% !TeX encoding = UTF-8
\chapter{Discussion and Conclusion}\label{chap:discussion}

\paragraph{Results.} All eight models detect the DDoS flows of this dataset very well (F1 from 0.9909 to 1.0000). The tree ensembles and KNN are best, with 0--6 errors on 33,463 test flows (XGBoost 0, Random Forest 3, KNN 6); their order is within the noise stated in the limitations. KNN has no training cost but the slowest prediction (4.204\,s for the test set). Among the neural networks the MLP is best (F1 0.9995), followed by the LSTM (0.9992) and the 1D-CNN (0.9991); none of them beats the tree ensembles, which also train in seconds instead of minutes. The LSTM needed tuning: with the shared settings it trained unstably, and gradient clipping, a lower learning rate and a longer patience improved it on every test metric. Gaussian Naive Bayes is the weakest model (F1 0.9909) because of 327 false alarms, followed by Logistic Regression (44 false positives). For a detector, the false alarms and missed attacks decide which model to prefer, more than the accuracy does (Figure~\ref{fig:errors}).

\paragraph{From detection to enforcement.} The models of Tasks~1--3 only classify flows, and our XDP filter (Chapter~\ref{ch:task4}) only limits packet rates. Connecting them needs a block list in the XDP program that a user-space detector fills (\cref{sec:t4-bridge}); since this happens once per flow, the prediction times of \cref{sec:cost} would be acceptable. We did not build this connection, and the rate limiter works without any model.

\paragraph{Limitations.} (1) The destination port is a shortcut: almost all DDoS flows target port 80, so a model can score high without learning the attack behavior, and we did not run an ablation. (2) We use one dataset from one day, one attack type and one split; the differences between models of a few errors are within what another seed or split could change. (3) The validation scores of the early-stopped models (XGBoost, MLP, LSTM, 1D-CNN) are optimistic. (4) The first choice of the epoch limits of the MLP and the CNN was made by comparing test metrics; we later re-justified the limits on validation evidence, and the effect was small (MLP F1 0.9994 $\to$ 0.9995; 1D-CNN F1 0.9986 $\to$ 0.9991, FP 21 $\to$ 13, FN 32 $\to$ 20), but the test set is therefore not perfectly untouched for these two choices. The final LSTM configuration was selected by validation loss; the test scores of the earlier LSTM variants were known when we decided to tune further. (5) Hyper-parameters are fixed and not searched.

\paragraph{Reproducibility.} All seeds are fixed. Results are reproducible in the same software environment; in a different environment the neural-network results can differ in the last decimals (a few errors), and the LSTM is the most sensitive (Chapter~\ref{chap:task2}). Such differences should not be over-interpreted.

\paragraph{Conclusion.} For this dataset, tree ensembles (XGBoost, Random Forest) and KNN are the best, and the tree ensembles are also cheap,, and neural networks give no advantage. Since the data is easy and partly shortcut-driven, harder data, other days and other attack types are needed before drawing general conclusions.
